\chapter{量子材料实验与数据阅读}

\begin{prereqbox}
理解倒空间、态密度、费米面、线性响应和基本相变。本章不训练仪器操作，而训练“测量量是什么、怎样从数据到物理结论、替代解释是什么”。
\end{prereqbox}
\begin{mapbox}
\[
\begin{array}{c|c}
\text{问题}&\text{常用探针}\\\hline
\text{原子/磁结构}&X\text{ 射线、中子、电子衍射}\\
\text{占据能带 }E(\vb k)&\text{ARPES}\\
\text{局域态密度}&\text{STM/STS}\\
\text{费米面与散射}&\text{输运、量子振荡}\\
\text{体相能隙与激发}&\text{比热、磁化率、光学、NMR}
\end{array}
\]
\end{mapbox}

\section{衍射：从峰位和强度到结构}
弹性散射满足
$\vb Q=\vb k_f-\vb k_i=\vb G$。峰位给晶格周期和对称性，强度
\[
I(\vb G)\propto
\left|\sum_j f_j(\vb G)e^{-i\vb G\cdot\vb r_j}\right|^2
\]
给基元信息。有限晶粒尺寸导致峰展宽，热振动给 Debye--Waller 衰减，无序产生弥散散射。

X 射线主要看电子密度，重元素更强；中子核散射长度不随原子序数单调，能看轻元素和同位素，且中子磁矩可测磁结构；电子散射强，适合薄样品和局域结构但多重散射更显著。

\begin{examplebox}
若降温后在原晶格禁戒位置出现新峰，先判断其波矢 $\vb Q$：整数/半整数超结构峰可提示晶格倍周期、磁序或电荷序。必须比较 X 射线与中子、是否有偏振分析，以及峰强随温度的序参量行为，不能仅凭“多了一个峰”断言具体序。
\end{examplebox}

\section{非弹性散射：测量动力学结构因子}
能量和动量同时转移时测得
\[
S(\vb Q,\omega)=\frac1{2\pi N}
\int\dd t\,e^{i\omega t}
\langle\rho_{\vb Q}(t)\rho_{-\vb Q}(0)\rangle.
\]
它是密度或自旋关联函数的频谱。峰的位置给激发色散，宽度给寿命，积分强度给谱权重。声子、磁振子和连续粒子--空穴激发可由其 $\vb Q$ 依赖、偏振和温度行为区分。

\section{ARPES：能量与动量分辨的占据谱}
光电效应给
\[
E_{\rm kin}=h\nu-\phi-E_B.
\]
平行表面的晶体动量近似守恒：
\[
k_\parallel=\frac{\sqrt{2mE_{\rm kin}}}{\hbar}\sin\theta.
\]
在突然近似和矩阵元框架下
\[
I(\vb k,\omega)\propto |M_{fi}|^2
A(\vb k,\omega)f(\omega).
\]
$A$ 是单粒子谱函数，$f$ 说明 ARPES 主要看占据态。峰位置给重整化色散，峰宽与自能虚部和寿命相关，强度还受光子能量、偏振、轨道选择定则与表面条件影响。

\begin{pitfallbox}
ARPES 不是“直接拍到裸能带”。实验强度包含矩阵元、分辨率、背景和 Fermi 函数；强关联时峰可能不再对应长寿命准粒子。表面态与体态可通过光子能量导致的 $k_z$ 色散、样品终止面和计算比较来区分。
\end{pitfallbox}

\section{STM/STS：实空间与局域谱}
弱隧穿、针尖态密度平坦时，Tersoff--Hamann 图像给
\[
I(V)\propto\int_0^{eV}\rho_s(\vb r,E_F+\epsilon)\dd\epsilon,
\qquad
\frac{\dd I}{\dd V}\propto\rho_s(\vb r,E_F+eV).
\]
恒流形貌是几何高度与积分态密度的混合，不是纯粹原子高度。STS 的
$\dd I/\dd V$ 近似局域态密度，可看超导相干峰、能隙、杂质束缚态和 Landau 能级。

准粒子干涉中，缺陷把 $\vb k$ 散射到 $\vb k'$，实空间态密度出现
$\vb q=\vb k'-\vb k$ 调制；Fourier 变换可反推等能轮廓和选择定则。缺失某散射矢量可能来自自旋纹理或矩阵元，但也可能是信噪比和无序。

\section{直流输运：电阻、Hall 与磁阻}
金属通常降温电阻下降，绝缘行为常表现为 $\dd\rho/\dd T<0$，但这不是严格相定义。超导转变要看零电阻、抗磁和热力学证据；仅有电阻下降可能是丝状通道。

Hall 斜率在单带模型给载流子符号和密度，多带时需拟合
$\rho_{xy}(B)$ 与纵向磁阻。反常 Hall 可写
\[
\rho_{xy}=R_0B+R_sM,
\]
第二项与磁化、Berry 曲率和散射机制有关。非线性 Hall 不自动证明多带，也可能含磁性或几何混合。

\section{量子振荡}
Shubnikov--de Haas 测电阻振荡，de Haas--van Alphen 测磁化振荡。对 $1/B$ 作 Fourier 变换得到频率和费米面极值面积。Lifshitz--Kosevich 温度因子
\[
R_T=\frac{X}{\sinh X},\qquad
X=\frac{2\pi^2\kb Tm^*}{\hbar eB}
\]
可拟合回旋质量；场依赖衰减给 Dingle 温度和量子寿命。Landau fan 截距与 Berry 相位有关，但 Zeeman、三维曲率和索引约定都会改变相位判断。

\section{比热、磁化率和热导}
连续相变常在比热出现跳变或 $\lambda$ 异常。超导弱耦合 BCS 有
$\Delta C/\gamma T_c\simeq1.43$；偏离可来自强耦合、多带、节点或样品体积分数。

零场冷却与场冷却磁化可看屏蔽和磁通钉扎。Curie 尾可能来自少量顺磁杂质，不应直接解释为本征磁性。热导由电子、声子和其他中性激发共同承担，结合 Wiedemann--Franz 与场依赖可分离通道。

\section{光学、Raman 与泵浦探测}
反射率经 Kramers--Kronig 可得复电导。Drude 峰看自由载流子，有限频率峰可来自带间跃迁、声子、激子或关联激发。总谱权重满足光学和规则；谱权重跨较大能标转移是关联的重要线索。

Raman 散射近似零动量转移，可按偏振选择晶格对称通道；泵浦--探测先把系统推离平衡，再追踪序参量、准粒子和集体模的恢复时间。非平衡信号不一定等于平衡能带简单加热。

\section{四类读图模板}
\begin{examplebox}
\textbf{判断超导。}电阻到仪器零值；低场磁化出现体抗磁；比热在同一 $T_c$ 有异常；隧穿或光学见能隙。若只有其中一项，应主动列出接触、丝状相、磁杂质或结构转变等替代解释。
\end{examplebox}
\begin{examplebox}
\textbf{判断磁相变。}磁化率出现 cusp 或发散、比热异常、中子出现磁 Bragg 峰。铁磁还看自发磁化和磁滞；反铁磁可能净磁矩为零，不能因磁化小而排除。
\end{examplebox}
\begin{examplebox}
\textbf{判断金属--绝缘跨越。}结合 $\rho(T)$、Hall、光学 Drude 权重、ARPES 是否有费米面及低温极限。局域化、能隙、颗粒效应和接触问题都可产生上翘电阻。
\end{examplebox}
\begin{examplebox}
\textbf{判断拓扑表面态。}ARPES 看跨越体隙的表面色散和自旋纹理，STS 看缺陷散射与 Landau 能级，输运看二维通道和量子化响应。任何单一“Dirac 锥”都不足以替代拓扑不变量与对称性分析。
\end{examplebox}

\begin{selftestbox}
\begin{enumerate}
\item 衍射峰位、强度、宽度分别告诉什么？
\item X 射线和中子对磁结构的能力有何不同？
\item $S(\vb Q,\omega)$ 与关联函数有何联系？
\item ARPES 强度为何不等于谱函数本身？
\item STM 形貌为什么不是纯几何高度？
\item 单带 Hall 分析何时失效？
\item 量子振荡怎样得到有效质量？
\item 怎样用三类证据确认体超导？
\item 比热、磁化率、输运对相变各有什么盲区？
\item 面对一张“漂亮数据图”，如何系统列出替代解释？
\end{enumerate}
\end{selftestbox}

